\section*{Supplementary Methods}
\setcounter{subsection}{0}
\renewcommand{\thesubsection}{SM\arabic{subsection}}
Detailed implementation and provenance accompany the main Methods. Internal evaluation always denotes the validation set used for checkpoint selection; later external analyses reuse a historically examined collection.

\subsection{WSI preprocessing and feature extraction}

WSI preprocessing used TRIDENT, with GrandQC tissue segmentation \citep{ref12}. Acquisition parameters were recovered from the archived WSI headers and coordinate-file attributes. All 193 internal and 168 analyzed external WSIs recorded Hamamatsu C9600-12 as the scanner model, a \(40\times\) source-lens setting and a level-0 resolution of approximately \(0.225708\,\mu\mathrm{m}\) per pixel in both axes. The original TIFF resolution tags were 44,305 pixels per centimetre. Internal files recorded NDP.scan versions 2.5.85 (133 WSIs) or 3.3.0 (60 WSIs), and all analyzed external files recorded version 2.5.85. These are file-recorded acquisition parameters; clinical provenance and stain identities are addressed separately. Supplementary Table S7 summarizes the recovered metadata.

Tissue regions were tiled at target magnification \(20\times\) into non-overlapping 512 \(\times\) 512-pixel patches. The archived HDF5 attributes consistently recorded a corresponding 1,024 \(\times\) 1,024-pixel level-0 footprint and zero overlap. Patch features and coordinates were stored in HDF5 files.

The principal experiments used UNI2-h embeddings with 1,536 dimensions per patch. The UNI family provides general-purpose pathological representations \citep{ref8}; the UNI2-h model card identifies the specific encoder family \citep{ref29}. The archived encoder implementation resized the shorter image side to 224 pixels, applied a 224 \(\times\) 224 centre crop, converted images to tensors and normalized RGB channels with means (0.485, 0.456, 0.406) and standard deviations (0.229, 0.224, 0.225). Its configuration specified a ViT-giant encoder with 14-pixel internal token patches, 24 transformer layers, 24 attention heads, eight register tokens and bfloat16 inference precision. The 14-pixel token size is distinct from the WSI extraction patch size. \missing{[MISSING M01: recover the exact UNI2-h weight-file checksum and downloaded checkpoint revision used for feature extraction.]} The image encoder was not fine-tuned within the downstream MIL experiment because training used precomputed embeddings. Each WSI contributed 512 sampled feature vectors. Sampling was without replacement when sufficient instances were available and with replacement otherwise. Training sampling varied across loading events, while inference sampling used deterministic seeds.

CONCH v1.5 with TITAN was evaluated as a separate comparison \citep{ref13}. It used 768-dimensional patch features and 1,024 sampled instances per WSI. Thus, its encoder, patch budget and aggregation strategy differed from those of the UNI2-h comparisons. The corresponding result is identified separately in main Figure 2 and Table 2.

% Restored from ../methods_hybrid_evidence_mil.tex and checked against
% ../../scripts/ss_experiment_suite.py: GatedAttentionPool,
% HybridEvidenceMILEncoder and SlideClassifier. This file is maintained by hand.
\subsection{Hybrid-Evidence MIL}

\paragraph{Task formulation and feature projection.}
Each WSI was represented by a sampled feature bag
$X=\{\mathbf{x}_i\}_{i=1}^{N}$ and its associated patch coordinates
$(a_i,b_i)$, where $N=512$ and $\mathbf{x}_i\in\mathbb{R}^{1536}$.
The slide label was $y=1$ for MALT and $y=0$ for SS.
No patch-level labels were used. A shared projection placed all features in a
$D=256$ dimensional hidden space:
\begin{equation}\label{eq:projection}
 \mathbf{h}_i = \operatorname{Dr}_{0.3}\!\left[
  \operatorname{GELU}\!\left(\operatorname{LN}
  (\mathbf{W}_{p}\mathbf{x}_i+\mathbf{b}_{p})\right)\right],
 \qquad \mathbf{h}_i\in\mathbb{R}^{D}.
\end{equation}
Here $\operatorname{LN}$ denotes layer normalization and $\operatorname{Dr}_{0.3}$
denotes dropout with probability 0.3 during training; dropout is disabled during
inference. All vectors below are column vectors, and $[\,\cdot\,;\,\cdot\,]$
denotes concatenation along the feature dimension.

\paragraph{Gated attention evidence.}
Gated pooling followed the attention mechanism of Ilse and colleagues
\citep{ref5}, with an additional value transformation implemented in the archived
model. For a sequence $Z=\{\mathbf{z}_j\}_{j=1}^{L}$ of $D$ dimensional vectors,
define the pooling operator $\mathcal{G}_{\theta}$ as
\begin{subequations}\label{eq:gated-pool}
\begin{align}
 \mathbf{t}_j &= \operatorname{Dr}_{0.3}
     \!\left[\operatorname{ReLU}(\mathbf{W}_{v}\mathbf{z}_j+\mathbf{b}_{v})\right],
     \label{eq:pool-value}\\
 \ell_j &= \mathbf{w}_{a}^{\mathsf T}
    \left[\tanh(\mathbf{V}\mathbf{t}_j+\mathbf{b}_{V})
     \odot\operatorname{sigmoid}(\mathbf{U}\mathbf{t}_j+\mathbf{b}_{U})\right]+b_a,
     \label{eq:pool-logit}\\
 \alpha_j &= \frac{\exp(\ell_j)}{\sum_{l=1}^{L}\exp(\ell_l)},
 \qquad \mathcal{G}_{\theta}(Z)=\sum_{j=1}^{L}\alpha_j\mathbf{t}_j.
     \label{eq:pool-sum}
\end{align}
\end{subequations}
The symbol $\odot$ denotes element-wise multiplication. All attention and value
matrices map between $D$ dimensional spaces. The attention evidence vector was
$\mathbf{e}_{\mathrm{att}}=\mathcal{G}_{\theta_A}(\{\mathbf{h}_i\}_{i=1}^{N})$.
Thus, the pooling weights multiply transformed values $\mathbf{t}_i$.
The regional branch used the same operator with independently learned parameters
$\theta_R$.

\paragraph{Distribution evidence.}
The center and heterogeneity of the sampled bag were summarized by two separate
vectors:
\begin{subequations}\label{eq:distribution}
\begin{align}
 \mathbf{e}_{\mu} &= \frac{1}{N}\sum_{i=1}^{N}\mathbf{h}_i,\\
 \mathbf{e}_{\sigma} &=
 \left[\frac{1}{N-1}\sum_{i=1}^{N}
       (\mathbf{h}_i-\mathbf{e}_{\mu})^{\odot 2}\right]^{\odot 1/2}.
\end{align}
\end{subequations}
Powers and square roots in Eq.~\eqref{eq:distribution} are element-wise. The
sample denominator $N-1$ matches the standard deviation used in the code.

\paragraph{Prototype-margin evidence.}
Two latent prototype banks indexed by $c\in\{0,1\}$ each contained
$K=8$ learnable prototypes
$\mathbf{p}_{c,r}\in\mathbb{R}^{D}$. Let a hat denote the implemented
$\ell_2$ normalization, with its numerical safeguard for zero norms.
The maximum cosine similarity within each bank and the bank-1-minus-bank-0 margin
were
\begin{subequations}\label{eq:prototype}
\begin{align}
 s_{i,c} &= \max_{1\leq r\leq K}
       \widehat{\mathbf{h}}_i^{\mathsf T}\widehat{\mathbf{p}}_{c,r},
       \qquad m_i=s_{i,1}-s_{i,0},\label{eq:margin}\\
 k &= \min(32,N),\qquad
 I_{+}=\operatorname{TopKIdx}(\{m_i\},k),\qquad
 I_{-}=\operatorname{TopKIdx}(\{-m_i\},k),\\
 \mathbf{e}_{+} &= \frac{1}{k}\sum_{i\in I_{+}}\mathbf{h}_i,
 \qquad \mathbf{e}_{-}=\frac{1}{k}\sum_{i\in I_{-}}\mathbf{h}_i.
\end{align}
\end{subequations}
$\operatorname{TopKIdx}$ returns indices of the $k$ largest entries. The high-
and low-margin vectors therefore average the instances with the largest and
smallest margins, respectively; the selected values need not all have the same
sign. Prototypes were optimized through differentiable similarity and margin
statistics in the slide-level classification objective. The hard top-$k$ index
selections were discrete; no additional prototype-specific loss or regional
labels were used. The bank indices do not identify validated MALT or SS
morphology, and a larger margin is not constrained to increase the final MALT score.

\paragraph{Coordinate-ordered groups and low-margin evidence.}
To form coordinate-ordered groups, let $\pi$ be the permutation obtained by
sorting the implemented coordinate key $10^6b_i+a_i$. The ordered hidden
sequence was partitioned using the code's chunk operation into $Q\leq16$
consecutive non-empty chunks $C_j$ (16 equal chunks for the 512 sampled patches).
Group tokens and the combined low-margin/regional vector were
\begin{subequations}\label{eq:regional}
\begin{align}
 \mathbf{r}_j &= \frac{1}{|C_j|}\sum_{l\in C_j}\mathbf{h}_{\pi(l)},
        \qquad j=1,\ldots,Q,\\
 \mathbf{e}_{\mathrm{reg}} &=
       \mathcal{G}_{\theta_R}(\{\mathbf{r}_j\}_{j=1}^{Q}),
       \qquad \mathbf{e}_{\mathrm{nr}}=\mathbf{e}_{-}+\mathbf{e}_{\mathrm{reg}}.
\end{align}
\end{subequations}
The regional operator includes its own value transformation from
Eq.~\eqref{eq:pool-value}. Coordinates determine group membership, but the
subsequent pooling is invariant to the order of the group tokens. The branch
has no explicit intergroup distances, adjacency or positional encoding.
Equal-sized groups of sampled instances were not anatomical regions and can
span different physical extents.

\paragraph{Similarity and margin statistics.}
Eight scalar statistics supplemented the evidence vectors. Define
$T_c=\operatorname{TopKIdx}(\{s_{i,c}\},k)$,
$\bar{s}_{c}^{\mathrm{top}}=k^{-1}\sum_{i\in T_c}s_{i,c}$ and
$\bar{m}=N^{-1}\sum_i m_i$. Then
\begin{subequations}\label{eq:statistics}
\begin{align}
 s_m &= \left[\frac{1}{N-1}\sum_{i=1}^{N}(m_i-\bar m)^2\right]^{1/2},\\
 \rho_0 &= \frac{1}{N}\sum_{i=1}^{N}\mathbb{I}(m_i>0),\qquad
 \rho_{\bar m}=\frac{1}{N}\sum_{i=1}^{N}\mathbb{I}(m_i>\bar m),\\
 \mathbf{u} &= \bigl[\bar{s}_{0}^{\mathrm{top}},\bar{s}_{1}^{\mathrm{top}},
       \bar m,s_m,\max_i m_i,\min_i m_i,\rho_0,\rho_{\bar m}\bigr]^{\mathsf T}
       \in\mathbb{R}^{8}.
\end{align}
\end{subequations}
$\mathbb{I}$ is the indicator function. The top-similarity sets $T_c$ are computed
separately for the two banks and need not coincide with the margin-based sets
$I_{+}$ and $I_{-}$.

\paragraph{Learned gating, weighted concatenation and classification.}
The five evidence vectors were ordered as
$(\mathbf{e}_1,\ldots,\mathbf{e}_5)=
(\mathbf{e}_{\mathrm{att}},\mathbf{e}_{\mu},\mathbf{e}_{\sigma},
 \mathbf{e}_{+},\mathbf{e}_{\mathrm{nr}})$.
A two-layer gate received the unweighted vectors and statistics:
\begin{subequations}\label{eq:fusion}
\begin{align}
 \mathbf{q} &= [\mathbf{e}_1;\mathbf{e}_2;\mathbf{e}_3;\mathbf{e}_4;
                   \mathbf{e}_5;\mathbf{u}]\in\mathbb{R}^{5D+8},\\
 \boldsymbol{\gamma} &= \operatorname{softmax}\!\left(
    \mathbf{W}_{g2}\operatorname{Dr}_{0.3}\!\left[
       \operatorname{GELU}(\mathbf{W}_{g1}\mathbf{q}+\mathbf{b}_{g1})\right]
       +\mathbf{b}_{g2}\right)\in\mathbb{R}^{5},\label{eq:gate}\\
 \widetilde{\mathbf{q}} &= [\gamma_1\mathbf{e}_1;\gamma_2\mathbf{e}_2;
       \gamma_3\mathbf{e}_3;\gamma_4\mathbf{e}_4;\gamma_5\mathbf{e}_5;
       \mathbf{u}],\label{eq:weighted-concat}\\
 \mathbf{z} &= \operatorname{Dr}_{0.3}\!\left[
       \operatorname{GELU}\!\left(\operatorname{LN}
       (\mathbf{W}_{f}\widetilde{\mathbf{q}}+\mathbf{b}_{f})\right)\right],\\
 a &= \mathbf{w}_{o}^{\mathsf T}\operatorname{Dr}_{0.3}\!\left[
       \operatorname{ReLU}\!\left(\operatorname{LN}
       (\mathbf{W}_{c}\mathbf{z}+\mathbf{b}_{c})\right)\right]+b_o,
       \qquad p=\operatorname{sigmoid}(a).
\end{align}
\end{subequations}
The gate dimensions were $\mathbf{W}_{g1}\in\mathbb{R}^{D\times(5D+8)}$
and $\mathbf{W}_{g2}\in\mathbb{R}^{5\times D}$;
$\mathbf{W}_{f}\in\mathbb{R}^{D\times(5D+8)}$ and
$\mathbf{W}_{c}\in\mathbb{R}^{D\times D}$.
Each $\gamma_j$ is non-negative and $\sum_j\gamma_j=1$.
Equation~\eqref{eq:weighted-concat} preserves five weighted feature blocks;
the eight statistics enter fusion without multiplication by a branch weight.
The gate supplies sample-dependent feature scaling, rather than a distribution
of causal contributions. Branch norms, the learned fusion transformation and
the ungated statistics also affect the classifier. The output $p$ is an
uncalibrated MALT model score in $[0,1]$.

\paragraph{Training objective.}
For a minibatch $\mathcal{B}$ of slides, the weighted binary cross-entropy was
\begin{equation}\label{eq:loss}
 \mathcal{L}=-\frac{1}{|\mathcal{B}|}\sum_{b\in\mathcal{B}}
 \left[w_{+}y_b\log p_b+(1-y_b)\log(1-p_b)\right],
 \qquad w_{+}=\frac{77}{67}.
\end{equation}
The implementation evaluated this objective directly from logits for numerical
stability. The positive-class weight was the ratio of the 77 SS to 67 MALT
training slides. Class weighting changes the optimum of binary cross-entropy;
the sigmoid output was not assumed to equal an unweighted posterior probability.
No calibration model was fitted to either evaluation set. Optimization and the
archived checkpoint are described in the main Methods.


\subsection{Fixed-checkpoint comparison and repeated input sampling}

The fixed-checkpoint analyses reused the four archived Hybrid, ABMIL, CLAM-SB-style and RRTMIL-style checkpoints without updating their weights or extracting new image features. A recorded protocol specified the original 49 internal and 168 external evaluation slides, a common threshold of 0.5 and the sampling analyses before inspection of the new outputs. Reference bags were reconstructed using the historical row indices and sampling rules. In particular, the original full external-manifest indices were retained when evaluating the final subset. All four internal outputs and the archived external Hybrid outputs reproduced to a maximum absolute score difference of $1.20\times10^{-7}$, with identical classifications at 0.5. The three comparators were then evaluated on exactly the same external feature bags as Hybrid.

Sampling sensitivity was measured with 20 additional 512-instance bags for every evaluation WSI, using the fixed Hybrid checkpoint in evaluation mode. For each slide and repeat 1--20, the NumPy random-number seed was the unsigned 32-bit integer represented by the first four bytes, in little-endian order, of the SHA-256 digest of \texttt{20261104|\allowbreak{}scope|\allowbreak{}study\_slide\_id|\allowbreak{}repeat\_id}. Sampling used replacement only when fewer than 512 feature rows were available. As in the historical implementations, indices were sorted only for internal slides with at least 512 rows. Exact indices, seeds, feature-row counts and sampled unique-row counts were saved.

For each slide we calculated the sample standard deviation, minimum and maximum of the 20 scores. A threshold-crossing slide had at least one score below 0.5 and at least one at or above 0.5 among those draws. Complete-cohort AUROC was calculated separately for each repeat. The distribution across draws measures technical input-sampling variation; it is not a confidence interval, an estimate of training variability or an increase in biological sample size. Averaging the 20 scores per slide was evaluated descriptively as a secondary result, without replacing the archived primary predictions.

On each reference bag, a separate instrumented forward calculation recorded latent-bank margins, prototype winner counts, gate scalings and vector norms. Its logits agreed with the original forward calculation within $10^{-5}$. Spearman correlations related the bag-mean margin to the output score and binary label. No regional annotations or causal interpretations were assigned to these internal quantities.

External complete-model AUROC differences used 10,000 class-stratified paired WSI bootstrap draws with seed 20261104 and fixed reference inputs. Each draw sampled the same slides for all models. Percentile intervals and exact paired McNemar tests were descriptive and unadjusted for multiple comparisons. Supplementary Table S11 uses this new 10,000-draw procedure, while the original primary Hybrid intervals are retained from the historical 2,000-draw analysis.


\subsection{Pathologist comparison and interpretation materials}

The available reader material comprised aggregate confusion counts for two pathologists. Individual decisions and the actual review conditions are held by the hospital data provider and await incorporation. The available aggregate counts are retained as provisional summaries pending reconciliation with those records. \missing{[MISSING R01: hospital-held individual reader decisions and protocol, including experience, randomization, blinding and case presentation; pending incorporation.]} Paired performance analyses remain unfilled pending record import. Reading time, confidence and assisted outcomes will only be reported if supported by the original study records.

The reader-analysis specification uses the same WSI-level diagnostic metric definitions and Wilson proportion intervals as the model analysis. Once matched reader records are available, reader--model error overlap and exact McNemar comparisons can be calculated on the same slides. These reader comparisons have not been performed from the provisional aggregates. Confidence-score ROC analysis requires the original reader scores; binary reader decisions alone will be reported through their corresponding diagnostic operating point.

Interpretation used saved case scores, gate weights, prototype-margin overlays and ranked tissue thumbnails. Four displayed cases represent a true positive, true negative, false positive and false negative at threshold 0.5. The performance and interpretation scripts each sampled 512 patches, using seed 42 plus a row index. Performance sampling indexed the evaluation subset; interpretation sampling indexed the full cohort table. Consequently, the sampled bags could differ. Supplementary Table S6 links both archived scores; the original figure outputs were linked without regenerating them.

The margin overlays retain their original pixels and depict sampled patch locations. High and low values describe relative similarity to two latent prototype banks, without assigning either color a verified disease morphology. Colors used a symmetric range $[-s,s]$, where $s=\max(|Q_{0.02}(m)|,|Q_{0.98}(m)|,10^{-6})$ within each slide. The figure legend displays normalized margin $m/s$, with color saturation beyond $\pm1$. Color intensity is therefore not directly comparable across slides. The six continuous similarity/margin statistics provide differentiable paths for prototype optimization. The two indicator proportions and hard top-$k$ indices are discrete operations.

Photomicrographs and ranked tissue thumbnails were reused without contrast, brightness or local image enhancement. Each thumbnail originated from a 1,024 $\times$ 1,024-pixel level-0 window spanning approximately 231.1 $\times$ 231.1 $\mu\mathrm{m}$. This calibration supports 50-$\mu\mathrm{m}$ bars on the model-selected thumbnails, which were stored at 220 $\times$ 220 pixels. For all eight displayed windows, the original 1,024-pixel RGB region resized to 220 pixels by Lanczos resampling matched the archived thumbnail pixel for pixel. Locator boxes use the archived overlay geometry. For Case 10, the archived coordinate metadata recorded a slide height of 62,208 pixels, whereas the current OpenSlide reader reported 70,400 pixels; both widths were 102,400 pixels. The original overlay and its coordinate geometry were retained rather than treating the current reader extent as a newly verified match. This calibration does not apply to the separate illustrative photomicrographs. \missing{[MISSING P01/P03: specimen linkage and original calibration for the illustrative photomicrographs, and pathology-confirmed regional descriptions.]}

Automated segmentation-contour coverage was audited for the same eight windows (Supplementary Table S12). Coverage was the area of the intersection between the union of archived segmentation polygons and the level-0 1,024 $\times$ 1,024-pixel window, divided by the window area. The original segmentation used $1\times$ magnification and treated polygon holes as tissue. This geometric overlap is an automated-mask measure, not a pathologist-confirmed tissue fraction. No ROI was removed or replaced on the basis of the audit; regional pathological interpretation remains pending hospital record incorporation.


\subsection{Component analyses}

Five configurations compared combinations of evidence vectors within the Hybrid architecture (Supplementary Table S9). Each used up to 30 epochs, batch size 16, learning rate $5\times10^{-5}$ and seed 42. When fewer than four vectors were active, all eight similarity/margin statistics were zeroed. Thus, the high-margin configuration simultaneously enabled that vector and the eight statistics. The final vector combined low-margin instances and coordinate-group summaries, without separating their contributions. Inactive vectors were masked, but learned gates retained a softmax over five output positions. The attention-vector configuration therefore retained the Hybrid gate and fusion architecture and was not identical to ABMIL. In the full equal-gate configuration, all five vectors were multiplied by 0.2 before concatenation. A fixed nonzero scaling can be absorbed into the subsequent linear fusion weights in the represented function class, although the parameterization can still affect optimization and regularization.

The seed controlled partitioning and the initialized global random-number streams. However, the component-training dataset created an unseeded NumPy \texttt{default\_rng()} for each training patch sample. It therefore did not fully fix patch-sampling randomness across runs. These results describe the saved configurations, rather than controlled attribution of independent components or exactly reproducible training trajectories.

The archived seed analysis contained seeds 7, 13, 21 and 42. The first three runs allowed 20 epochs and seed 42 allowed 30 epochs. We reported individual AUROCs without treating the four runs as interchangeable repetitions. These historical runs are shown in Supplementary Fig. S3. Ensemble behavior was described by aligned score correlations and by the individual errors corrected, introduced or shared at threshold 0.5.


\subsection{Controlled comparison across training seeds}

A separate protocol specified three configurations and training seeds 17, 29 and 43, for nine runs. Matched attention used the same 1,536-to-256 projection, gated attention block and 256-input classification head as Hybrid, while omitting the other vectors, statistics and fusion layer. Fixed scaling retained all five vectors, the eight statistics and the fusion layer, with each vector multiplied by 0.2 and no trainable gate. Learned gating used the full original Hybrid architecture. Shared tensors started from identical freshly initialized values within each seed. These configurations were not parameter-matched, and the fixed-versus-learned gate comparison also changed nonlinear capacity.

All runs used the existing 144/49 split, 30 epochs, batch size eight, AdamW with learning rate $5\times10^{-5}$ and weight decay $10^{-4}$, dropout 0.3 and positive BCE weight 77/67. The first epoch achieving a strictly higher internal AUROC than all preceding epochs was retained, and final-epoch checkpoints were also archived. External predictions used the selected checkpoint and were computed after selection. There was no early stopping, learning-rate scheduler or calibration fitting.

For each seed and epoch, the configurations received identical sampled training bags and minibatch orders. Bag seeds were the first four little-endian SHA-256 bytes of \texttt{controlled-v6|\allowbreak{}seed|\allowbreak{}epoch|\allowbreak{}study\_slide\_id}; order seeds replaced the slide code with \texttt{order}. PyTorch batch seeds used \texttt{controlled-v6|\allowbreak{}seed|\allowbreak{}epoch|\allowbreak{}batch\_index|\allowbreak{}dropout}. These reset common random streams, although dropout masks after divergent computational branches were not assumed to be elementwise matched. Exact indices and minibatch orders were saved, deterministic algorithms were enabled, and parameter-gradient connectivity was checked in the first minibatch of each run. Internal and external evaluation used the same fixed reference bags as the checkpoint analyses above.

The nine predefined runs were reported individually, with configuration means and sample standard deviations and paired per-seed AUROC differences (Supplementary Table S15). The three seeds describe training variation on the fixed cohort. They were not used as independent clinical observations or to establish inferential significance. The original ABMIL implementation and all historical checkpoint results remained separate.


\subsection{Statistical analysis and figure generation}

The primary display threshold was 0.5. Archived optimized thresholds were reported separately and were not relabeled as prespecified. AUROC measured rank discrimination, and AUPRC was calculated as average precision \citep{ref18}. Accuracy, sensitivity, specificity, PPV, NPV and F1 were calculated from the corresponding confusion matrix. Balanced accuracy was the mean of sensitivity and specificity. The constant-positive accuracy reference equalled the observed MALT fraction. The unit of these analyses was the WSI.

Metrics were computed from saved model scores using the accompanying Python scripts. The original 95\% intervals for AUROC, average precision and F1 used 2,000 label-stratified bootstrap samples with fixed predictions. The random seed was 20261016. Proportion intervals used the Wilson method \citep{ref19}. These intervals condition on the saved score for each WSI and do not incorporate variation from resampling the input bag. No unprovided patient linkage was invented. Hybrid-versus-comparator AUROC differences used paired resampling of aligned predictions. Exact McNemar tests used discordant classifications at the same threshold. These fixed-score intervals also omit checkpoint/model selection and unknown patient clustering. The original two pairwise internal model comparisons were descriptive and were not multiplicity-adjusted. Brier scores were the mean squared differences between the uncalibrated scores and binary labels. Their 95\% percentile intervals used 10,000 class-stratified WSI bootstrap samples with fixed archived scores; model-specific seeds 20261104--20261109 are archived in the source table. This describes squared prediction error in each cohort, rather than isolating calibration or validating clinical probabilities.

The retained v7 quantitative figures were generated in Python with PDF/SVG exports and accompanying source tables. Figure 1 retains its vector cohort panel and uses an AI-assisted model schematic checked against the implementation. Performance-comparison tables retain the archived values and mark the best and second-best distinct point estimates within the stated comparison groups. Tied displayed values receive the same marking; confidence limits, cohort counts and non-comparable groups are not ranked. These descriptive ranks do not indicate statistical significance. Histology panels retain source tissue pixels; all additional annotations are separate graphical elements. Outstanding legacy-image calibration and pathology annotations are marked. The current figure and statistical environment used Python 3.10.18, NumPy 1.26.4, pandas 2.3.2, SciPy 1.15.3, scikit-learn 1.7.2 and Matplotlib 3.10.6. New checkpoint inference used PyTorch 2.8.0 with CUDA 12.8, float32 and a single NVIDIA A800 80GB GPU. The fixed-checkpoint inference run took 43.23 seconds, excluding environment setup and subsequent statistical analysis; it did not extract image features or update model weights. The separate nine-run training comparison took 185.44 seconds including loading the cached features and checkpoint evaluation. Both timings concern the current machine and cached-feature workloads. This current environment is recorded separately from the historical training environment.

Implementation provenance was documented separately from this statistical environment. The archived TRIDENT package identifies version 0.2.0; its checkout revision begins a6305acfef68. The custom experiment script and saved principal Hybrid checkpoint are identified by SHA-256 prefixes 5de055ea93926ec5 and ccdffb739d5b65ee, respectively. Full hashes, the saved configuration, encoder-source identity and recovered metadata accompany the draft. These identify the inspected archived files and do not substitute for the original runtime environment. \missing{[MISSING M02: recover the original training GPU and historical Python, PyTorch, CUDA and dependency versions.]}

\subsection{Random-state records for the primary analysis}

The original label-stratified 144/49 split used random seed 42 and a holdout fraction of 0.25. The primary Hybrid training run also used seed 42. Its highest internal validation AUROC selected epoch 15 from the 30-epoch history. External outcomes did not select or replace this checkpoint. Bootstrap random states and deterministic evaluation-bag conventions are specified above and in the accompanying analysis code.
